\section{粒子与空穴算符}

Gross 第 19 章把 Wick 定理建立在相对费米海的粒子/空穴算符上。这比“对任意真空收缩”更贴近固体物理的语言。

\subsection{定义}

取无扰基态为填满的费米海 $|\Phi_0\rangle$，单粒子能级以 $\varepsilon_F$ 为界。定义
\begin{equation}
\begin{aligned}
  a_{\mathbf{k}}&=c_{\mathbf{k}},
  \quad
  a_{\mathbf{k}}^\dagger=c_{\mathbf{k}}^\dagger
  && (k>k_F),\\
  b_{\mathbf{k}}&=c_{\mathbf{k}}^\dagger,
  \quad
  b_{\mathbf{k}}^\dagger=c_{\mathbf{k}}
  && (k<k_F).
\end{aligned}
\label{eq:ph_ops}
\end{equation}
前者为粒子算符，后者为空穴算符。
则
\begin{equation}
  a_{\mathbf{k}}|\Phi_0\rangle=b_{\mathbf{k}}|\Phi_0\rangle=0,
\end{equation}
即 $|\Phi_0\rangle$ 是粒子--空穴真空。$a$ 与 $b$ 相互反对易（作用在不同轨道）。

\subsection{自由哈密顿量}

\begin{equation}
  H_0
  =W_0
  +\sum_{k>k_F}\xi_{\mathbf{k}}a_{\mathbf{k}}^\dagger a_{\mathbf{k}}
  +\sum_{k<k_F}|\xi_{\mathbf{k}}|b_{\mathbf{k}}^\dagger b_{\mathbf{k}},
\end{equation}
$\xi_{\mathbf{k}}=\varepsilon_{\mathbf{k}}-\mu$。激发能 $=$ 粒子能量之和 $+$ 空穴能量之和（空穴能量为 $\mu-\varepsilon_{\mathbf{k}}>0$）。

\subsection{场算符的分解}

\begin{equation}
  \psi(\mathbf{r})
  =\psi^{(+)}(\mathbf{r})+\psi^{(-)}(\mathbf{r}),
\end{equation}
$\psi^{(+)}$ 含空穴产生（$k<k_F$ 的 $b^\dagger$），$\psi^{(-)}$ 含粒子湮灭（$k>k_F$ 的 $a$）。相互作用绘景中
\begin{equation}
  a_{\mathbf{k}}(t)=a_{\mathbf{k}}e^{-i\xi_{\mathbf{k}}t},
  \qquad
  b_{\mathbf{k}}(t)=b_{\mathbf{k}}e^{-i|\xi_{\mathbf{k}}|t}.
\end{equation}

\subsection{正规序与收缩}

粒子/空穴算符乘积的正规序 $N(\cdots)$ 把所有 $a^\dagger,b^\dagger$ 移到左边。收缩定义为
\begin{equation}
  \overline{A B}
  =T(AB)-N(AB),
\end{equation}
非零收缩只有
\begin{equation}
  \overline{a_{\mathbf{k}}(t)a_{\mathbf{p}}^\dagger(t')}
  =\delta_{\mathbf{k}\mathbf{p}}\theta(t-t')e^{-i\xi_{\mathbf{k}}(t-t')},
\end{equation}
以及空穴的 $\theta(t'-t)$ 项。这正是自由 $G_0$ 的粒子/空穴分解。Wick 定理：编序乘积 $=$ 全部完全收缩 $+$ 含正规序的部分；对真空期望只留完全收缩。

\begin{note}
图中向下的空穴线与向上的粒子线，即来自此分解。等时线必须规定 $t'\to t^+$ 以免占据数歧义。
\end{note}
